\section{The finite law and its support}\label{sec:law}

Write
\[
 \PhiPot(S)=\sum_{v\in S}v,
 \qquad G_S(z)=\E[z^{T_S}].
\]

\begin{proposition}[absorption and PGF]\label{prop:pgf}
Every trajectory is absorbed after at most $\PhiPot(S)$ updates.  Moreover,
\begin{equation}
 G_{\{0\}}(z)=1,
 \qquad
 G_S(z)=\frac{z}{|S|-1}
       \sum_{v\in S\setminus\{0\}}G_{C_v(S)}(z).       \label{eq:pgf}
\end{equation}
Thus \eqref{eq:pgf} is a finite acyclic recursion for the complete law.
\end{proposition}

\begin{proof}
If $v-1\notin S$, the update decreases $\PhiPot$ by one.  If
$v-1\in S$, the collision removes $v$ and decreases $\PhiPot$ by $v$.
The nonnegative integer potential therefore decreases strictly.  Conditioning
on the uniformly chosen first active site gives \eqref{eq:pgf}; potential
descent makes the recursion acyclic.
\end{proof}

The recursion has no internal gaps, even for a sparse initial state.

\begin{theorem}[support from every rooted state]\label{thm:support}
If $S\ne\{0\}$, then
\begin{equation}
 \supp(T_S)=\{\max S,\max S+1,\ldots,\PhiPot(S)\}.       \label{eq:support}
\end{equation}
\end{theorem}

\begin{proof}
Each update decreases $\PhiPot$ by at least one, so $T_S\leq\PhiPot(S)$.
An update lowers the maximum occupied coordinate by at most one: moving the
maximum either places it at its predecessor or removes it when that
predecessor is occupied; a move below the maximum leaves it fixed.  Hence
$T_S\geq\max S$.

We prove that every intermediate value occurs by induction on
$\PhiPot(S)$.  Put $m=\max S$.  If $m-1\notin S$, choose $m$ first.  The
successor has maximum $m-1$ and potential $\PhiPot(S)-1$, so induction and
the first update realize the whole interval $[m,\PhiPot(S)]$.

Suppose instead that $m-1\in S$.  The case $m=1$ is immediate.  Choosing
$m$ first is a collision; induction supplies
\begin{equation}
 [m,\PhiPot(S)-m+1].                                   \label{eq:lowinterval}
\end{equation}
Let $a$ be the bottom of the consecutive occupied run ending at $m$.
If $a=1$, moving $a$ collides at the root; if $a>1$, its predecessor is
empty.  Either move decreases the potential by one and leaves maximum $m$.
This first choice and induction supply
\begin{equation}
 [m+1,\PhiPot(S)].                                     \label{eq:highinterval}
\end{equation}
Since $m,m-1\in S$, we have $\PhiPot(S)\geq2m-1$; the integer intervals in
\eqref{eq:lowinterval}--\eqref{eq:highinterval} touch or overlap.  Their
union proves \eqref{eq:support}.  Every constructed schedule has positive
probability because each of its active choices does.
\end{proof}

Let $S_n=\{0,1,\ldots,n-1\}$ and $T_n=T_{S_n}$.

\begin{corollary}[minimum-time mass]\label{cor:minmass}
For $n\geq2$,
\begin{equation}
 \Pp(T_n=n-1)=\frac{1}{(n-1)!}.                         \label{eq:minmass}
\end{equation}
\end{corollary}

\begin{proof}
A trajectory of length $n-1$ must reduce the number of nonroot piles at
every update, so every move must collide.  Removing a site before its larger
neighbor would leave a vacancy for that neighbor's next move.  The unique
all-collision order is therefore $n-1,n-2,\ldots,1$.  The successive numbers
of active sites are $n-1,n-2,\ldots,1$, giving \eqref{eq:minmass}.
\end{proof}

